\documentclass[11pt]{article}

\usepackage[margin=1.1in]{geometry}
\usepackage{amsmath,amssymb,amsthm}
\usepackage{mathtools}
\usepackage{stmaryrd}
\usepackage{microtype}
\usepackage[T1]{fontenc}
\usepackage{palatino}
\usepackage{enumitem}
\usepackage[hidelinks]{hyperref}

% ---- theorem environments ----
\theoremstyle{definition}
\newtheorem{definition}{Definition}
\newtheorem{principle}{Principle}
\theoremstyle{plain}
\newtheorem{proposition}{Proposition}
\newtheorem{conjecture}{Conjecture}
\newtheorem{question}{Question}
\theoremstyle{remark}
\newtheorem{remark}{Remark}
\newtheorem*{obligation*}{Obligation}
\newtheorem{obligation}{Obligation}

% ---- process-calculus notation ----
\newcommand{\quoteP}{@}                       % quote
\newcommand{\deref}{{*}}                       % dereference
\newcommand{\inn}[3]{\mathrm{for}(#1 \leftarrow #2)\,#3}
\newcommand{\out}[2]{#1\,!\,(#2)}
\newcommand{\nil}{\mathbf{0}}
\newcommand{\rcong}{\equiv}
\newcommand{\bisim}{\sim}
\newcommand{\rbisim}{\sim_{r}}
\newcommand{\reduces}{\longrightarrow}
\newcommand{\sem}[1]{\llbracket #1 \rrbracket}
\newcommand{\bagunion}{\uplus}
\newcommand{\Mfin}{M_{\mathrm{fin}}}
\newcommand{\Name}{\mathrm{Name}}
\newcommand{\Proc}{\mathrm{Proc}}
\newcommand{\dgr}{\mathbf{d}}
\newcommand{\Cut}{\mid}

\title{\bfseries Not Every Cut Is an Aperture\\[2pt]
\large Choice as a computational resource, confluent interiors,\\
and the seams along which agents interfere}

\author{Lucius Gregory Meredith\\
\small CEO, F1R3FLY.io, and Chief Science Officer, F1R3FLY Industries\\
\small \texttt{lgreg.meredith@gmail.com}}

\date{2026}

\begin{document}
\maketitle

\begin{abstract}
The mobile process calculi (MPCs) furnish a vocabulary in which several
ontological questions about computational agency become sharp. Working
within that vocabulary --- and treating the RSpace semantics of the
reflective higher-order calculus as a \emph{motivating example} rather than
as a model of any physical mechanism --- we examine when one computational
agent can interfere with another. Our organising observation is that the
nondeterminism of an interaction is a property of the \emph{cut} at which two
processes meet, not of either process taken alone, and that the resolution of
that nondeterminism is a computational resource: a choice function for the
family of admissible pairings, graded in strength by exactly the infinitude
the source calculus's risk ledger already charges. An agent of higher
computational capacity can influence a weaker one only by writing into such a
resource where it is left open --- but not every cut leaves it open. We
isolate the \emph{macro-component}: a sub-term whose internal interaction is
confluent up to resource bisimilarity, so that however much its interior
branches, no internal cut is observable as choice; its only apertures lie at
its boundary. Eric Smith's account of the biosphere --- the agents that
consume energy from outside as the true interface to the rest of the universe,
the agents that consume one another as a rigid interior --- is, in this
vocabulary, exactly such a factorisation. Linearity is too strong a criterion
for the rigid interior; confluence is the right one, and the $\lambda$-calculus,
operationally awash in branching yet famously determinate, is the paradigm.
We close on the one structural question on which the stability of the whole
picture turns: whether confluence up to resource bisimilarity is preserved
under subspace descent --- equivalently, whether boundary interference can
crack open a previously rigid interior, which would be the formal content of
\emph{perturbation}.
\end{abstract}

\section{Introduction}

This note is a conceptual essay carried by a notation. It is about agency,
interference, and the boundaries between computational entities of unequal
capacity; and its method is to borrow the syntax and semantics of the mobile
process calculi (MPCs) --- the $\pi$-calculus and, specifically, the reflective
higher-order (rho) calculus \cite{rho} --- not as a physical model but as an
\emph{ontological vocabulary}: a language precise enough to name the categories
at play and to say which distinctions are real and which are artefacts of how we
choose to describe.

Two disclaimers fix the register. First, the RSpace denotation of the rho
calculus \cite{qcs,paths} appears throughout as a \emph{motivating example}. We
are not positing that RSpace, its keyed multiset store, or its ``forgiving zip''
bears any resemblance to a physical mechanism. We invoke it because it makes one
thing vivid --- that the indeterminacy of which sender meets which receiver is
localised at a cut and is, formally, a choice --- and because that one thing is
enough to seed the rest. Second, where we reach for biology we lean on Smith and
Morowitz's \emph{The Origin and Nature of Life on Earth} \cite{smith}, whose
reading of the biosphere as a planetary process with a definite interface to the
rest of the universe is, we will argue, illuminated by --- and illuminating of
--- the MPC vocabulary. The claim is not that life \emph{is} a process term. The
claim is that the categories Smith draws (the agent that eats outside energy; the
agent that eats other agents; the biosphere-as-agent versus the bee-as-agent) are
categories the MPC vocabulary already has names for, and that putting the two
side by side sharpens both.

The thesis the note exists to state is a single sentence, and it is about cuts.
In the MPCs, two agents always sit on either side of a parallel composition; the
composition is the locus of their interaction, the \emph{cut}. One is tempted to
read every such cut as an opening --- a place where what one agent does is
underdetermined and therefore a place where another agent can reach in. The
temptation is wrong, and correcting it is the work of the note:

\begin{principle}[Not every cut is an aperture]\label{prin:main}
Agents flank every cut. But only some cuts leave the interaction's
nondeterminism \emph{open and observable}, and only at those cuts can an agent
of higher computational capacity influence the behaviour of a lower one.
The cuts that do are the \emph{apertures}; the cuts that do not are sealed by
the determinacy --- not necessarily the linearity, but the confluence --- of
the agents that meet there.
\end{principle}

\noindent
The body develops, in order: choice as a resource (\S\ref{sec:choice}); the
correction that locates the aperture at the cut rather than in the agent
(\S\ref{sec:aperture}); the resulting picture of a flat ``jumble'' of unequal
agents versus a tower of oracles, and the conditions under which the jumble is
coherent (\S\ref{sec:jumble}); Smith's biosphere as a factorisation in the
vocabulary (\S\ref{sec:smith}); the macro-component, with confluence (not
linearity) as the criterion for a rigid interior, and the $\lambda$-calculus as
the witness (\S\ref{sec:macro}); the coincidence of three frontiers and what it
buys (\S\ref{sec:frontier}); and the one open structural question on which the
stability of the whole construction rests (\S\ref{sec:crux}).

\paragraph{The MPC vocabulary, briefly.}
We assume the modern presentation of \cite{qcs}. Processes and names are given
by
\[
P,Q ::= \nil \mid P \Cut Q \mid \inn{y}{x}{P} \mid \out{x}{Q} \mid \deref x,
\qquad x,y ::= \quoteP P,
\]
with $\nil$ the stopped process, $P\Cut Q$ parallel composition (the cut),
$\inn{y}{x}{P}$ input on channel $x$ binding $y$, $\out{x}{Q}$ asynchronous
output, $\deref x$ dereference, and a name solely a quoted process $\quoteP P$
(no primitive names, no restriction $\nu$). Structural congruence $\rcong$ makes
$(\,\Cut\,,\nil)$ a commutative monoid and validates the reflection law
$\deref\quoteP P \rcong P$. The base reaction is communication at a cut,
\[
\inn{y}{x}{P} \Cut \out{z}{Q} \reduces P\{\quoteP Q / y\}
\qquad (x \equiv_N z),
\]
and the two equivalences we use are context bisimilarity $\bisim$ (the
behavioural equivalence, a congruence by construction when transitions are
labelled by contexts \cite{leifermilner,sewell,oslf}) and resource
bisimilarity $\rbisim$ (the finer relation that counts parallel copies). The
RSpace denotation $\sem{-}\colon \Proc \to T(\Proc)$ interprets $\Cut$ as
keyed multiset union $\bagunion$, so $\sem{-}$ is a monoid homomorphism and
communication is the ``zip'' of dual bags at a key, with the surplus left as
residue \cite{qcs}. None of what follows depends on the set-theoretic model of
\cite{qcs,paths}; it depends only on the shape of the vocabulary.

\section{Choice as a computational resource}\label{sec:choice}

Begin with the one feature of the RSpace example we will actually use. When a
channel carries both a bag of consumes $\{\!|(y_1)P_1,\dots,(y_n)P_n|\!\}$ and a
bag of produces $\{\!|Q_1,\dots,Q_m|\!\}$, communication chooses an underlying
list of each and zips them, keeping unmatched payloads as surplus. The choice of
list orderings is the nondeterminism of the calculus --- which send meets which
receive --- now localised at the key. The denotation itself remains a function;
reaction is the separate dynamics; and the indeterminacy is exactly the choice of
pairing.

That is not a metaphor for choice. It \emph{is} a choice function, for a specific
family. Index by the cuts that carry mixed polarity (or, for a run, by the
reaction events); at each such index the fibre is the set of admissible pairings,
inhabited precisely when a redex is present. A completed schedule is a section of
this family --- an element of the product of the fibres --- and the assertion that
such a section exists is an instance of the axiom of choice. The scheduler is the
choice function for the pairing family, not by analogy but by construction. This
gives the practical reframing we want: \emph{choice is a computational resource},
the resource of nondeterminism, and different ways of supplying it import
different amounts of computational power.

\paragraph{The resource is graded by the same $\omega$ the risk ledger charges.}
The strength of the choice principle one must invoke is not uniform; it tracks
the infinitude of the family.

\begin{itemize}[leftmargin=1.4em]
\item \emph{Finite bags, finite run.} The index set is finite, the product is
inhabited with no choice axiom at all (constructively, in $\mathrm{IZF}$). More
than that: a bag is a parallel composition of finitely many prefixes, so it
arrives \emph{presented} as a list-up-to-permutation. The term hands you the
enumeration; picking a linearisation is constructively free. The term is its own
choice structure.

\item \emph{One $\omega$.} Replication and reflection install an $\omega$-fold
bag, or a non-terminating run yields an $\omega$-sequence of events. Now one needs
countable choice for the static store, or --- for a run whose pairings depend on
the residuals of the previous step --- \emph{Dependent Choice}. The cascade of a
deep cut unpacking into a deeper cut \cite{paths} is exactly a dependent-choice
sequence. This is the single $\omega$ --- the ``black infinity'' --- the risk
ledger of \cite{qcs} already spends; choice-as-resource and the $\omega$-meter
are the same gauge.

\item \emph{Class-sized.} Full $\mathrm{AC}$, and only here.
\end{itemize}

\noindent
So the slogan ``choice is not non-constructive'' is exactly true in the finite
fragment and \emph{relocated, not dissolved}, above it. The constructive escape
at $\omega$ is to never form the completed schedule: take each pairing locally, on
demand, as a decidable step, replacing the choice \emph{function} (whose bare
existence needs $\mathrm{DC}$ or $\mathrm{AC}$) with a choice \emph{process}. That
is what lazy, on-demand reaction does operationally; it keeps choice potential
rather than actual. The honest reading is of a piece with the rest of the
programme: recursion relocated from anti-foundation into reflection, congruence
from an afterthought into the context labels, and now choice from a metatheoretic
existence axiom into an operationally consumed resource.

\paragraph{A lattice, not a tower.}
The natural home for ``choice as a resource'' already exists: the Weihrauch
lattice of Brattka, Gherardi and Pauly \cite{weihrauch}, whose objects are
choice principles ordered by uniform reducibility, with demonic binary choice
appearing as the lesser-limited-principle-of-omniscience and composition
$\star$ modelling ``use one resource, then another.'' The decisive structural
fact, for our purposes, is that the Weihrauch degrees form a \emph{lattice} with
rich algebra, not a linear hierarchy. We will return to this: the intuition that
a universe of computational agents is ``of all stripes, not neatly sorted by a
tower of oracles'' is, formally, the statement that the relevant degree structure
is a lattice and not an ordinal chain.

\section{The aperture is a property of the cut}\label{sec:aperture}

It is tempting to summarise \S\ref{sec:choice} by saying that an agent's
\emph{nondeterminism budget} is its interface to the rest of the universe: a
fully deterministic process is sealed, with every slot filled by its own
transition function and nowhere for anything external to write; a
nondeterministic agent is open precisely at its choice points. There is a true
idea here --- under-determination is the aperture --- but the locus is wrong, and
the parallel-composition algebra is what makes it wrong.

Nondeterminism is not preserved by $\Cut$ in either direction. A whole that is
deterministic can decompose into parts each wildly underdetermined at its
\emph{internal} cuts; and parts each internally rigid can compose into a whole
whose only indeterminacy is at the boundary between them. The aperture is
therefore a property of the \emph{cut}, not of either agent at it, and --- this is
the consequence that matters --- \emph{which} cut one privileges as ``the agent
boundary'' is a modelling choice. Different observers draw it at different depths
of the same term, and the interior lights up or goes dark accordingly.

\begin{remark}[Under-determination versus boundary]\label{rmk:seam}
Apertures coincide with under-determination: a cut is an aperture exactly when the
resolution of its pairing nondeterminism is \emph{observable}, i.e.\ when distinct
admissible pairings yield residuals in distinct $\rbisim$-classes (equivalently,
distinct $\bisim$-classes, after the projection that forgets multiplicity). But
the seam at which under-determination sits need not respect any particular
delineation of agents. A seam may lie far above, or far below, the boundary an
observer would draw; its effects may be felt at cuts the observer counts as
internal. This is why ``the agent's budget is its interface'' fails: the budget is
real, but it belongs to the cut, and the cut may straddle whatever boundary one
had in mind.
\end{remark}

In the path reading of \cite{paths}, where keys are routes through a trie and the
prefix order gives one channel a subspace of another, the modelling freedom
becomes literal: \emph{a choice of agent boundary is a choice of prefix}, and the
trie is the lattice of possible boundaries. A coarse boundary is a short prefix,
near the root, folding everything below into one sub-store and exposing only a
shallow aperture; a fine boundary is a long prefix, exposing apertures at every
internal branch. These are not rival claims about where the boundary ``really''
is. They are different truncation depths on one structure, ordered by the prefix
order, which is precisely the order ``coarser agent / finer agent.'' The In-deep
rule of \cite{paths} --- a reaction at a shallow cut unpacking a cut that lives
deeper than the boundary-drawer was tracking --- is what it looks like when the
same dynamics is read at two depths at once.

\section{Stratification versus the jumble}\label{sec:jumble}

The natural ontology suggested by ``choice imports computational power'' is a
\emph{tower of oracles}: agents stratified by degree, the stronger able to reach
down into the weaker. We want to register a different intuition and then say
exactly what saves it. A universe made of computational phenomena may be not a
neat tower but a \emph{jumble} --- agents of all stripes and capabilities
interacting, as we see in the physical world, with the stratification only a
framework that aids understanding, not the structure of the world. If that is the
case, agents with strong oracles could interfere with the ``choices'' of agents
less computationally endowed, but they do so as participants in a flat field, not
as occupants of a higher floor.

\paragraph{Both pictures at once: a (bi)fibration.}
The two readings need not compete. Let a total category $\mathcal{E}$ be the
jumble --- every agent of every degree, with the cut-mediated maps between them ---
and let a projection to the degrees (Turing, Weihrauch, or an ordinal index) be
the stratification. Reindexing along $\dgr' \le \dgr$ --- the cartesian, fibration
direction --- takes a degree-$\dgr$ computation to its degree-$\dgr'$ shadow, the
view a weaker observer keeps after forgetting the oracle. That is the
``stratification as aid to understanding'' move made precise: \emph{the tower is a
cleavage, a chosen gauge for describing a frame-independent jumble}. Interference
is the other direction, and for it one wants the reindexing functors to carry
Lawvere adjoints $\exists_f \dashv f^{*} \dashv \forall_f$. We record, as a
conjecture rather than a theorem, the fit that motivates the formalisation:

\begin{conjecture}[Angelic/demonic as adjoints]\label{conj:adjoint}
Angelic resolution (the cooperative scheduler that \emph{finds} a good pairing)
and demonic resolution (the adversarial scheduler that must survive \emph{every}
pairing) are, respectively, the left and right Lawvere adjoints
$\exists_f \dashv f^{*} \dashv \forall_f$ to reindexing along the
degree/prefix order. If so, the scheduler ladder --- canonical, demonic-finite,
angelic, fair-merge, relativised, transfinite --- is the adjoint structure of one
fibration rather than a list of cases.
\end{conjecture}

\noindent
The load-bearing check is whether the fibres carry the limits and colimits to
support both adjoints; that is where a formalisation should begin.

\paragraph{What keeps the jumble from collapsing.}
A flat ontology of \emph{unequal} agents is only coherent if a single strong agent
cannot simply dominate. Two design decisions of the source vocabulary are exactly
the consistency conditions. \emph{No implicit interaction} \cite{paths}: the store
never reacts with itself; a reaction is triggered only by a cut that actually
co-locates opposite polarities. Drop it, and a single halting-oracle agent on a
universal channel homogenises the field to the top degree --- the jumble dies.
\emph{Locality of the cut}: interference can be written only where a channel is
shared, and, in the path reading, only into the subspace below a held prefix.
Position in the trie and degree \emph{together} set the reach of interference. Add
a cost discipline (the phlogiston of the source programme) and reach becomes finite
per unit resource. Locality plus cost is what keeps a multi-degree universe a
jumble rather than a hierarchy that flattens upward. We state this as a slogan,
because it is the most decision-relevant thing the vocabulary offers: \emph{no
implicit interaction and the risk ledger are not bookkeeping; they are the
consistency conditions for a flat ontology of unequal computational agents.}

\section{Smith's biosphere as a factorisation}\label{sec:smith}

Smith and Morowitz \cite{smith} read terrestrial life as a planetary process and
ask after its interface to the rest of the universe. In their account the agents
that matter at that interface are the ones that consume energy from \emph{outside}
the biosphere --- sunlight, the chemical disequilibria of hydrothermal vents ---
while a great interior of agents consume \emph{one another}, trading already-captured
energy. The energy-eaters are, on this reading, the actual boundary between
(terrestrial) life and the cosmos: at the level of the interface, the agentic
boundary makes the entire biosphere one agent. Yet humans, themselves life forms,
recognise a far finer delineation --- bees distinct from birds, this bee distinct
from that one.

The MPC vocabulary names these categories cleanly. Write, schematically,
\[
\textsf{Biosphere}
\;\rcong\;
\textsf{EnergyEaters} \;\Cut\; \textsf{AgentEaters},
\qquad
\textsf{Universe}
\;\rcong\;
\textsf{Biosphere} \;\Cut\; \textsf{RestOfUniverse}.
\]
The seam to the cosmos is the cut $\textsf{Biosphere} \Cut
\textsf{RestOfUniverse}$, and \textsf{EnergyEaters} is the sub-term whose
business is precisely that cut --- the place where outside energy, untyped to any
particular molecule (the sun does not address its photons to a chosen chlorophyll),
enters as an underdetermined produce. Meanwhile \textsf{AgentEaters}, taken as a
whole, may be entirely determinate. Crucially, this is so even though its internal
interactions are extremely constraint-sensitive: the bee consumes the pollen, not
the worm the bird eats; enzymes are spectacularly type-constrained, one substrate
per active site. The interior adheres to a thicket of conformance constraints and
yet has \emph{no races} --- the constraints do not select among matches, they
ensure there is only one match to select.

This is the factorisation Principle~\ref{prin:main} is about, and it inverts the
discrimination/choice relationship one might expect.

\begin{remark}[Two senses of ``agent'']\label{rmk:twoagents}
The bees are real agents to each other and to us --- distinct nodes, distinct
types, distinct rigid roles. But they need not be \emph{sites of choice}. Their
distinctness is type-distinctness, not aperture-distinctness. The only
underdetermined thing in a bee's life may be which boundary reaction, upstream,
gates the arrival of its pollen; given that arrival, everything is forced.
Discrimination among agents and location of nondeterminism have come apart. An
agent in the \emph{discrimination} sense (a typed, individuated role) is not, in
general, an agent in the \emph{choice} sense (a cut with a live race). Smith's
factorisation is the construction that separates the two: the finer, bee-level
description sees \emph{more agents} but, if the interior is rigid, the \emph{same
apertures} --- all sitting at the macro-boundary.
\end{remark}

\noindent
And it makes the seam-versus-boundary point of Remark~\ref{rmk:seam} concrete and
biological. The seam to the cosmos sits at the energy cut; its effects are felt at
the predation cuts, deep inside; and at that depth the agents may experience an
externally sourced bias as their own internal, determinate dynamics. No agent at
the predation level can, from inside its own capacity, distinguish ``my internal
choice resolved thus'' from ``the external write propagated down to here'' ---
\emph{provided} the interior is rigid enough to transmit the boundary decision
without adding choices of its own. When that proviso holds, the seam genuinely
violates the agent boundaries an observer would draw: it enters at one level and is
metabolised, at another, into what looks locally like determinism. Making the
proviso precise is the next section's task.

\section{The macro-component: confluence, not linearity}\label{sec:macro}

We want a name for the object Smith's interior instantiates: a sub-term all of
whose internal cuts are sealed, with apertures only at its boundary. The first
guess is that the interior is sealed by \emph{rigidity} --- that at every internal
key the consume-bag and produce-bag admit a \emph{unique} compatible pairing, so
that the forgiving zip is a function and there is no race to resolve. Determinism
not by scheduling but by there being nothing to schedule. Syntactically this is a
\emph{linearity} condition: every internal channel is used at multiplicity one per
polarity, each consume a linear hypothesis discharged by exactly one produce, the
boundary the residue of the non-linear, contended channels. On this reading
\textsf{EnergyEaters} are exactly the non-linear channels --- the ones contended by
many consumers, the sun feeding everyone --- and the factorisation is
\emph{linear interior, non-linear boundary}.

Linearity is, however, too strong. It is a syntactic sufficient condition for what
we actually care about, which is that the interior have no \emph{observable} choice.
The semantic statement of exactly that, and nothing more, is confluence.

\begin{definition}[Rigid sub-term / macro-component]\label{def:macro}
A sub-term is \emph{rigid} (internally) if its internal reaction is confluent up to
resource bisimilarity: every internal pairing-race reconverges, so all maximal
internal reaction sequences reach $\rbisim$-equal normal forms. A
\emph{macro-component} is a maximal rigid sub-term: every internal cut is a
non-aperture (confluent up to $\rbisim$), and its boundary cuts --- where
confluence fails --- are its apertures. The set of cuts at which internal
confluence fails is the \emph{rigidity frontier}.
\end{definition}

\noindent
Confluence permits the branching that linearity forbids. The bee may have several
internally admissible pairings; if they all flow to $\rbisim$-equal normal forms,
the bee is still not a site of observable choice, which is all the factorisation
needs. Metabolism is full of internal branching --- parallel pathways, redundant
enzymes, order-independent cascades --- that converges on the same steady state:
confluent, not linear. So the right slogan is \emph{the interior is
$\lambda$-like, not linear-like}.

\paragraph{$\lambda$ is really deterministic.}
The $\lambda$-calculus is the load-bearing witness, not an aside. Operationally it
is radically nondeterministic: a term with many redexes has many reduction
sequences, every interleaving a distinct branch. And it is \emph{the} paradigm of
determinism --- because Church--Rosser says the branching is fictional: one normal
form, reached however you like. The choice of redex is real and it is free in the
sense of \S\ref{sec:choice}, semantically free, spent but unobservable, eliminated
by the quotient to the normal form. $\beta$-reduction is the cleanest instance in
all of computing of a process saturated with races whose races do not matter, and
it is deterministic for precisely the reason Definition~\ref{def:macro} asks of an
interior: confluence, not absence of branching. A \emph{closed} $\lambda$-term ---
no free variables --- is then a sealed macro-component at every scale, nothing to
interfere with because there is no contended channel for anything to write to. An
\emph{open} term is the agent with an aperture: its free variables are its boundary
channels, the spots where a context (a higher agent) can substitute and change the
result. In the MPC vocabulary the free variable is the unmatched consume, the
dangling polarity --- the same object, the port where the universe reaches in.

\begin{proposition}[Degree is carried by the boundary]\label{prop:degree}
A confluent (up to $\rbisim$) interior contributes no choice resource: it carries
no Weihrauch content, no observable bits of nondeterminism, however baroque its
branching, because the resource was defined (\S\ref{sec:choice}) as observable
choice and confluence makes choice unobservable. Hence, modulo the obligations
below, the resource degree of a macro-component is carried entirely by its
boundary:
\[
\deg(\text{macro-component}) \;=\; \deg(\text{boundary}),
\]
the interior being degree-transparent. A linear-interior criterion would
under-count here, miscalling some genuinely confluent interiors ``boundary''
because they branch; the confluence criterion is what makes the equation clean.
\end{proposition}

\section{Three frontiers, one locus}\label{sec:frontier}

The reward of taking confluence as the criterion is that the boundary of the
macro-component coincides with two frontiers we already had names for.

\begin{proposition}[Frontier coincidence]\label{prop:frontier}
The following three loci coincide, for a macro-component in the sense of
Definition~\ref{def:macro}:
\begin{enumerate}[label=(\roman*),leftmargin=2.2em]
\item the \emph{rigidity frontier} --- the cuts where internal confluence up to
$\rbisim$ fails;
\item the \emph{detectable-interference frontier} --- the cuts where a biased
pairing yields a different $\rbisim$-class, so that an external agent's write
leaves an observable mark rather than being absorbed into a $\rbisim$-equal
outcome;
\item the \emph{confluence obligation frontier} --- the cuts at which the
eager-confluence-up-to-$\rbisim$ obligation of the path semantics
(\cite{paths}, Obligation~2) is non-trivial.
\end{enumerate}
\end{proposition}

\noindent
The coincidence is a consistency check on the vocabulary: the place where the
interior stops being able to hide its schedule is the place where external
interference starts being able to leave a mark, which is the place where the
formal confluence obligation bites. One frontier, three readings ---
confluence-failure from inside, interference-visibility from outside, proof
obligation from the meta-theory. That a single locus answers to all three is the
kind of agreement that suggests the apparatus is carving at a joint.

It also tells us when interference is a \emph{measurement}. Where internal reaction
is confluent up to $\rbisim$, a biased pairing picks the trajectory but not the
destination; the weak agent is steered to a $\rbisim$-equal normal form and cannot
tell. Where confluence fails, the bias becomes a different observable outcome.
Detectable interference is interference at a non-confluent cut --- an interaction
that resolves an otherwise-unobservable choice into an observable distinction,
which is close to the definition of a measurement. We note, without leaning on it,
that the way a weak agent could in principle detect a stronger one is
information-theoretic and contextual: its ``free'' choices would look random
relative to its own capacity yet be compressible relative to the higher degree,
carrying more mutual information with the oracle than chance allows --- local
sections admitting no degree-uniform global section, in the
sheaf-theoretic sense of contextuality \cite{abramsky}. The probabilistic face is
cleaner: a flat jumble is one joint distribution, an agent's view a marginal, and a
strong agent biasing a weak one's choices is conditioning that marginal on oracle
information the weak agent cannot access. We flag these as directions, not results.

\section{The crux: is rigidity stable under descent?}\label{sec:crux}

Everything above treats ``deterministic interior, nondeterministic boundary'' as a
static property of a term. Whether it is also \emph{stable under the dynamics} is a
single structural question, and it is the question on which the picture stands or
falls.

The danger is specific. Confluence up to $\rbisim$ is not compositional the way
linearity is: two confluent components composed at a boundary need not have a
confluent union --- that is the whole content of the boundary being where races
live. The subtler problem is the cascade. In the path semantics a deep cut
(In-deep) unpacks into a \emph{fresh, deeper} cut \cite{paths}; the confluence of
that manufactured cut is a new obligation not discharged by the confluence of the
parts. So the rigid-interior property defined by confluence is not obviously closed
under the very reaction --- subspace descent --- that propagates a boundary
decision inward.

\begin{question}[Stability of rigidity]\label{q:crux}
Is confluence up to resource bisimilarity preserved under subspace descent (the
In-deep cascade)? Equivalently: if a macro-component's boundary aperture fires and
the decision cascades inward, does the interior remain rigid?
\end{question}

\noindent
This is Obligation~2 of \cite{paths} wearing the clothes of the agent
factorisation, and the two possible answers are both informative.

\emph{If rigidity is stable.} The factorisation is robust: a boundary race fires,
the decision descends as a \emph{forced} cascade through an interior that has no
opinion of its own (only whether its unique-up-to-$\rbisim$ match is present), and
the macro-component survives as a structure that can be reasoned about across the
dynamics. ``Deterministic interior, nondeterministic boundary'' is then a genuine
structure theorem, and the degree equation of Proposition~\ref{prop:degree} holds
not just statically but along reductions. This is the outcome Smith's biology
quietly assumes: the predation web metabolises the day's sunlight into ordinary,
determinate dynamics.

\emph{If rigidity fails to be stable.} Then boundary interference can \emph{crack
open} a previously rigid interior --- turn a confluent agent into a branching one
by descending into it. We resist reading this as a mere defect. It is, plausibly,
the formal content of \emph{perturbation}: the higher agent does not merely bias the
interior's schedule among $\rbisim$-equal outcomes, it changes the interior's very
determinacy, opening apertures where there were none. On this reading a sufficiently
strong agent does not only steer a weaker one within its existing freedom; it can
\emph{manufacture} new freedom in the weaker one and then steer that. Which of these
two phenomena the vocabulary actually supports is decided by the answer to
Question~\ref{q:crux}, and the answer is checkable. It is the first thing a
formalisation should resolve, because it is the difference between a static fact
about terms and a dynamic, perturbable feature of them --- and, in the ontological
register of this note, the difference between a universe whose agents have fixed
capacities and one in which capacity itself is something agents do to each other.

\section{Obligations}

In the manner of \cite{qcs,paths}, the points a rigorous development would
discharge. Several are restatements, in the agent vocabulary, of obligations
already on the books there; that overlap is the point --- the conceptual claims of
this note reduce to structural claims about the source calculus.

\begin{obligation}[Choice as a graded resource]
Make precise the family-of-pairings whose section is the schedule, and prove the
grading of \S\ref{sec:choice}: finite runs need no choice; the single $\omega$ of
the risk ledger is exactly the step from no-choice to countable/dependent choice;
the In-deep cascade is a $\mathrm{DC}$ sequence; class-size forces full
$\mathrm{AC}$. Identify which schedulers are confluent up to $\rbisim$
(constructive, no oracle) and which are not (the imported oracle), and locate the
boundary at demonic-finite versus angelic-or-infinite.
\end{obligation}

\begin{obligation}[The adjoint structure]
Test Conjecture~\ref{conj:adjoint}: exhibit the (bi)fibration of agents over
degrees, and determine whether angelic and demonic resolution are the Lawvere
adjoints to reindexing. The decisive sub-question is whether the GSLT fibres carry
the limits and colimits both adjoints require.
\end{obligation}

\begin{obligation}[Macro-components are well defined]
Show that ``maximal rigid sub-term'' (Definition~\ref{def:macro}) is well defined:
that the rigidity frontier exists and is computable from term-plus-typing, and that
the three frontiers of Proposition~\ref{prop:frontier} coincide. Establish
Proposition~\ref{prop:degree} --- degree carried by the boundary --- rigorously,
including the claim that a confluent interior is resource-degree-transparent.
\end{obligation}

\begin{obligation}[Linearity implies, but is not implied by, rigidity]
Confirm that linear (multiplicity-one-per-polarity) internal typing implies
internal confluence up to $\rbisim$, and exhibit confluent-but-non-linear interiors
(redundant-pathway metabolism is the intended class) to show the implication is
strict --- so that confluence is the correct, weaker criterion.
\end{obligation}

\begin{obligation}[Stability of rigidity under descent --- the crux]
Resolve Question~\ref{q:crux}: decide whether confluence up to $\rbisim$ is
preserved by the In-deep cascade. If yes, prove the structure theorem
(deterministic interior / nondeterministic boundary, stable along reductions). If
no, characterise the perturbation --- the conditions under which a boundary write
opens new interior apertures --- as a property of the sub-trie.
\end{obligation}

\section{Conclusion}

We have used the mobile process calculi as a vocabulary, not a model. In that
vocabulary the indeterminacy of an interaction is a choice function for a definite
family, hence a computational resource, graded in strength by exactly the
infinitude the source calculus's risk ledger already charges; and an agent of
higher capacity can influence a weaker one only by writing into that resource where
it is left open. The correction the note turns on is that the opening is a property
of the \emph{cut}, not of either agent, so that ``which cut is the boundary'' is a
choice of description --- a choice of prefix, in the path reading --- and the same
seam may be an interface at one level and invisible internal dynamics at another.
Smith and Morowitz's biosphere, with its energy-eating interface and its
agent-eating interior, is exactly such a factorisation, and it sharpens into a
definite object: the macro-component, a sub-term whose interior is confluent up to
resource bisimilarity --- rigid not by linearity but in the way the
$\lambda$-calculus is deterministic, awash in branching that does not matter --- with
its apertures swept to a boundary that is simultaneously the rigidity frontier, the
detectable-interference frontier, and the source calculus's confluence obligation.
Two senses of agent come apart there: the individuated, typed role and the live
site of choice need not coincide, and Smith's interior shows they routinely do not.

The thesis is Principle~\ref{prin:main}: not every cut is an aperture. Agents flank
every cut; only some cuts leave the nondeterminism open and observable; only at
those can a stronger agent reach in. Whether that picture is static or dynamic ---
whether a boundary write merely steers a rigid interior or can crack it open and
manufacture new freedom within it --- is the single structural question of
Section~\ref{sec:crux}, and it is the same confluence-under-descent obligation the
path semantics already carries. We do not claim the universe is a process term. We
claim that when one asks where agents are, which of them can disturb which, and
along what seams, the MPC vocabulary names the categories well enough that the
questions become answerable --- and that the most consequential of them reduces to
a confluence lemma one can actually try to prove.

\paragraph{Acknowledgements.}
This note develops a line of discussion growing out of \cite{qcs,paths} and their
RSpace denotation, and out of Smith and Morowitz's reading of the biosphere
\cite{smith}; the conceptual claims here are deliberately routed back to structural
obligations on the source calculus, where they can be settled.

\begin{thebibliography}{99}

\bibitem{qcs} L.~G. Meredith. \emph{Quoting is Colour-Swap: a model of the rho
calculus in the knotted universe}. Manuscript, 2026.

\bibitem{paths} L.~G. Meredith. \emph{Paths are Subspaces: a cut-triggered
polymorphism of RSpace over path-keys, and its distributive law}. Manuscript, 2026.

\bibitem{knot} L.~G. Meredith. \emph{A Knotted Universe: a new notion of reflective
set theories}. Manuscript, 2026.

\bibitem{rho} L.~Gregory Meredith and Matthias Radestock. A reflective higher-order
calculus. \emph{Electr.\ Notes Theor.\ Comput.\ Sci.}, 141(5):49--67, 2005.

\bibitem{smith} Eric Smith and Harold J. Morowitz. \emph{The Origin and Nature of
Life on Earth: The Emergence of the Fourth Geosphere}. Cambridge University Press,
2016.

\bibitem{leifermilner} James J. Leifer and Robin Milner. Deriving bisimulation
congruences for reactive systems. In \emph{CONCUR 2000}, LNCS 1877, pages
243--258. Springer, 2000.

\bibitem{sewell} Peter Sewell. From rewrite rules to bisimulation congruences.
\emph{Theoret.\ Comput.\ Sci.}, 274(1--2):183--230, 2002.

\bibitem{oslf} L.~G. Meredith, Christian Wells, et al. Minimal-context transition
systems and the observer-specified logic functor for graph-structured lambda
theories. Manuscript, forthcoming.

\bibitem{turiplotkin} Daniele Turi and Gordon Plotkin. Towards a mathematical
operational semantics. In \emph{LICS 1997}, pages 280--291. IEEE, 1997.

\bibitem{rutten} J.~J.~M.~M. Rutten. Universal coalgebra: a theory of systems.
\emph{Theoret.\ Comput.\ Sci.}, 249(1):3--80, 2000.

\bibitem{weihrauch} Vasco Brattka, Guido Gherardi, and Arno Pauly. Weihrauch
complexity in computable analysis. In V. Brattka and P. Hertling, editors,
\emph{Handbook of Computability and Complexity in Analysis}, pages 367--417.
Springer, 2021.

\bibitem{abramsky} Samson Abramsky and Adam Brandenburger. The sheaf-theoretic
structure of non-locality and contextuality. \emph{New J.\ Phys.}, 13:113036, 2011.

\bibitem{aczel} Peter Aczel. \emph{Non-Well-Founded Sets}, volume 14 of
\emph{CSLI Lecture Notes}. CSLI Publications, Stanford, CA, 1988.

\end{thebibliography}

\end{document}
